\chapter{2005年浙江卷（理）}

\section{单选题}
\begin{problem}
极限 \(\displaystyle\lim_{n \to \infty} \frac{1+2+\cdots+n}{n^{2}}\) 的值是（\(\quad\)）

\choices
    {\(2\)}
    {\(4\)}
    {\(\frac{1}{2}\)}
    {\(0\)}

\begin{answer}
C
\end{answer}

\begin{solution}
由等差数列求和公式，\(1+2+\cdots+n=\dfrac{n(n+1)}{2}\)，所以
\[
\displaystyle\lim_{n\to\infty}\frac{1+2+\cdots+n}{n^{2}}
=\displaystyle\lim_{n\to\infty}\frac{n(n+1)}{2n^{2}}
=\displaystyle\lim_{n\to\infty}\frac{n+1}{2n}=\frac{1}{2}
\]
故选 C.
\end{solution}
\end{problem}

\begin{problem}
点 \((1,-1)\) 到直线 \(x-y+1=0\) 的距离是（\(\quad\)）

\choices
    {\(\frac{1}{2}\)}
    {\(\frac{3}{2}\)}
    {\(\frac{\sqrt{2}}{2}\)}
    {\(\frac{3\sqrt{2}}{2}\)}

\begin{answer}
D
\end{answer}

\begin{solution}
点 \((x_0,y_0)\) 到直线 \(Ax+By+C=0\) 的距离为
\(\dfrac{|Ax_0+By_0+C|}{\sqrt{A^2+B^2}}\). 因此所求距离为
\[
\frac{|1-(-1)+1|}{\sqrt{1^2+(-1)^2}}
=\frac{3}{\sqrt{2}}=\frac{3\sqrt{2}}{2}
\]
故选 D.
\end{solution}
\end{problem}

\begin{problem}
设 \( f(x)=\begin{cases}
|x-1|-2, & |x|\leqslant1,\\
\dfrac{1}{1+x^{2}}, & |x|>1,
\end{cases} \) 则 \( f\left[f\left(\frac{1}{2}\right)\right]= \)（\(\quad\)）

\choices
    {\( \frac{1}{2} \)}
    {\( \frac{4}{13} \)}
    {\(-\frac{9}{5}\)}
    {\( \frac{25}{41} \)}

\begin{answer}
B
\end{answer}

\begin{solution}
因为 \(\left|\frac{1}{2}\right|\leqslant1\)，所以应使用分段函数的第一段，得
\[
f\left(\frac{1}{2}\right)=\left|\frac{1}{2}-1\right|-2=-\frac{3}{2}
\]
又因为 \(\left|-\frac{3}{2}\right|>1\)，所以计算外层函数值时应使用第二段，得
\[
f\left[f\left(\frac{1}{2}\right)\right]
=f\left(-\frac{3}{2}\right)
=\frac{1}{1+\left(-\frac{3}{2}\right)^2}
=\frac{4}{13}
\]
故选 B.
\end{solution}
\end{problem}

\begin{problem}
在复平面内，复数 \( \frac{\i}{1+\i} + (1 + \sqrt{3}\i)^{2} \) 对应的点位于（\(\quad\)）

\choices
    {第一象限}
    {第二象限}
    {第三象限}
    {第四象限}

\begin{answer}
B
\end{answer}

\begin{solution}
先将分式分母实数化，得
\[
\frac{\i}{1+\i}=\frac{\i(1-\i)}{(1+\i)(1-\i)}=\frac{1+\i}{2}
\]
又
\[
(1+\sqrt{3}\i)^2=1+2\sqrt{3}\i+3\i^2=-2+2\sqrt{3}\i
\]
所以原复数为
\[
\frac{1+\i}{2}-2+2\sqrt{3}\i
=-\frac{3}{2}+\left(\frac{1}{2}+2\sqrt{3}\right)\i
\]
其实部为负，虚部为正，故对应的点位于第二象限，选 B.
\end{solution}
\end{problem}

\begin{problem}
在 \((1 - x)^{5} + (1 - x)^{6} + (1 - x)^{7} + (1 - x)^{8}\) 的展开式中，含 \(x^{3}\) 的项的系数是（\(\quad\)）

\choices
    {\(74\)}
    {\(121\)}
    {\(-74\)}
    {\(-121\)}

\begin{answer}
D
\end{answer}

\begin{solution}
由二项式定理，\((1-x)^j\) 中 \(x^3\) 的系数为
\(\mathrm C_j^3(-1)^3=-\mathrm C_j^3\). 因此原式中 \(x^3\) 的系数为
\[
-\left[\mathrm C_5^3+\mathrm C_6^3+\mathrm C_7^3+\mathrm C_8^3\right]
=-(10+20+35+56)=-121
\]
故选 D.
\end{solution}
\end{problem}

\begin{problem}
设 \(\alpha\)，\(\beta\) 为两个不同的平面，\(l\)，\(m\) 为两条不同的直线，且 \(l \subset \alpha\)，\(m \subset \beta\)，有如下的两个命题：
\begin{circlelist}
\item 若 \(\alpha \parallel \beta\)，则 \(l \parallel m\)；
\item 若 \(l \perp m\)，则 \(\alpha \perp \beta\).
\end{circlelist}
那么（\(\quad\)）

\choices
    {\(\circled{1}\)是真命题，\(\circled{2}\)是假命题}
    {\(\circled{1}\)是假命题，\(\circled{2}\)是真命题}
    {\(\circled{1}\circled{2}\)都是真命题}
    {\(\circled{1}\circled{2}\)都是假命题}

\begin{answer}
D
\end{answer}

\begin{solution}
命题\(\circled{1}\)不成立. 取 \(\alpha:z=0\)，\(\beta:z=1\)，并在其中分别取
\(l=\{(t,0,0)\mid t\in\R\}\)，\(m=\{(0,t,1)\mid t\in\R\}\).
此时 \(\alpha\parallel\beta\)，但两条直线的方向向量分别为 \((1,0,0)\) 和 \((0,1,0)\)，不平行，故命题\(\circled{1}\)是假命题.

命题\(\circled{2}\)也不成立. 取 \(\alpha:z=0\)，\(\beta:z=x\)，并取
\(l=\{(0,t,0)\mid t\in\R\}\)，\(m=\{(t,0,t)\mid t\in\R\}\).
两条直线都经过原点，且方向向量的内积为
\((0,1,0)\cdot(1,0,1)=0\)，所以 \(l\perp m\). 但两平面的法向量可分别取 \((0,0,1)\) 和 \((-1,0,1)\)，其内积为 \(1\)，两平面所成的锐二面角为 \(45^{\circ}\)，并非垂直. 故命题\(\circled{2}\)也是假命题，选 D.
\end{solution}
\end{problem}

\begin{problem}
设集合 \(A=\{(x,y)\mid x,y,1-x-y\text{ 是三角形的三边长}\}\)，则 \(A\) 所表示的平面区域（不含边界的阴影部分）是（\(\quad\)）

\choices
    {\choicebitmap[width=0.15\paperwidth]{img/2005/zhejiang_science/q7_optA_nolabel.png}}
    {\choicebitmap[width=0.15\paperwidth]{img/2005/zhejiang_science/q7_optB_nolabel.png}}
    {\choicebitmap[width=0.15\paperwidth]{img/2005/zhejiang_science/q7_optC_nolabel.png}}
    {\choicebitmap[width=0.15\paperwidth]{img/2005/zhejiang_science/q7_optD_nolabel.png}}

\begin{answer}
A
\end{answer}

\begin{solution}
令三边长分别为 \(x\)，\(y\)，\(1-x-y\). 它们能构成三角形的充要条件是三者均为正数，且每一边都小于另外两边之和. 因此
\(x>0\)，\(y>0\)，\(1-x-y>0\)，\(x<1-x\)，\(y<1-y\)，\(1-x-y<x+y\).
化简得
\(0<x<\frac{1}{2}\)，\(0<y<\frac{1}{2}\)，\(\frac{1}{2}<x+y<1\).
其中 \(x<\frac{1}{2}\) 和 \(y<\frac{1}{2}\) 已推出 \(x+y<1\)，所以区域是正方形
\(0<x<\frac{1}{2}\)，\(0<y<\frac{1}{2}\) 中直线 \(x+y=\frac{1}{2}\) 上方且不含边界的部分，对应图 A.
\end{solution}
\end{problem}

\begin{problem}
已知 \(k < -4\)，则函数 \(y = \cos 2x + k(\cos x - 1)\) 的最小值是（\(\quad\)）

\choices
    {\(1\)}
    {\(-1\)}
    {\(2k + 1\)}
    {\(-2k + 1\)}

\begin{answer}
A
\end{answer}

\begin{solution}
令 \(t=\cos x\)，则 \(-1\leqslant t\leqslant1\)，且
\[
y=2t^2-1+k(t-1)=2t^2+kt-k-1
\]
记右侧为 \(h(t)\)，则
\(h'(t)=4t+k\leqslant4+k<0\)，\((-1\leqslant t\leqslant1)\).
所以 \(h(t)\) 在 \([-1,1]\) 上单调递减，最小值在 \(t=1\) 处取得. 因而
\[
y_{\min}=2+k-k-1=1
\]
故选 A.
\end{solution}
\end{problem}

\begin{problem}
设 \( f(n) = 2n + 1\)（\(n \in \N\)），\( P = \{1, 2, 3, 4, 5\} \)，\( Q = \{3, 4, 5, 6, 7\} \)，记 \( \hat{P} = \{n \in \N \mid f(n) \in P\} \)，\( \hat{Q} = \{n \in \N \mid f(n) \in Q\} \)，则 \( \left(\hat{P} \cap \complement_{\N} \hat{Q}\right) \cup \left(\hat{Q} \cap \complement_{\N} \hat{P}\right) = \)（\(\quad\)）

\choices
    {\(\{0, 3\}\)}
    {\(\{1, 2\}\)}
    {\(\{3, 4, 5\}\)}
    {\(\{1, 2, 6, 7\}\)}

\begin{answer}
A
\end{answer}

\begin{solution}
由 \(f(n)=2n+1\) 以及 \(0\in\N\)，可得
\[
\hat P=\{n\in\N\mid 2n+1\in\{1,2,3,4,5\}\}=\{0,1,2\}
\]
\[
\hat Q=\{n\in\N\mid 2n+1\in\{3,4,5,6,7\}\}=\{1,2,3\}
\]
因此
\(\hat P\cap\complement_{\N}\hat Q=\{0\}\)，\(\hat Q\cap\complement_{\N}\hat P=\{3\}\)
从而所求集合为 \(\{0,3\}\)，选 A.
\end{solution}
\end{problem}

\begin{problem}
已知向量 \(\bs{a} \neq \bs{e}\)，\(|\bs{e}| = 1\)，对任意 \(t \in \R\)，恒有 \(|\bs{a} - t\bs{e}| \geqslant |\bs{a} - \bs{e}|\). 则（\(\quad\)）

\choices
    {\(\bs{a} \perp \bs{e}\)}
    {\(\bs{a} \perp (\bs{a} - \bs{e})\)}
    {\(\bs{e} \perp (\bs{a} - \bs{e})\)}
    {\((\bs{a} + \bs{e}) \perp (\bs{a} - \bs{e})\)}

\begin{answer}
C
\end{answer}

\begin{solution}
两边平方并利用 \(|\bs e|=1\)，得到对任意 \(t\in\R\)，
\[
|\bs a-t\bs e|^2-|\bs a-\bs e|^2
=t^2-2t\bs a\cdot\bs e+2\bs a\cdot\bs e-1
=(t-1)^2+2(t-1)(1-\bs a\cdot\bs e)\geqslant0
\]
令 \(t=1+h\)，则
\(h^2+2h(1-\bs a\cdot\bs e)\geqslant0\)，\((h\in\R)\).
若 \(1-\bs a\cdot\bs e\ne0\)，取充分小且符号相反的 \(h\)，上式将小于零，矛盾. 所以
\(\bs a\cdot\bs e=1\). 因而
\[
\bs e\cdot(\bs a-\bs e)=\bs a\cdot\bs e-|\bs e|^2=1-1=0
\]
又因 \(\bs a\ne\bs e\)，故 \(\bs e\perp(\bs a-\bs e)\)，选 C.
\end{solution}
\end{problem}

\section{填空题}
\begin{problem}
函数 \( y = \frac{x}{x + 2} \)（\(x \in \R\)，且 \(x \neq -2\)）的反函数是\fillinblank{}.

\begin{answer}
\(y=\frac{2}{1-x}\)，\(x\in\R\)，且 \(x\neq1\)
\end{answer}

\begin{solution}
由 \(y=\dfrac{x}{x+2}\) 得
\(y(x+2)=x\)，\(\Longrightarrow\)，\(x(y-1)=-2\).
由于原函数的值域为 \(\R\setminus\{1\}\)，所以 \(y\ne1\)，从而
\[
x=\frac{2}{1-y}
\]
交换 \(x\)，\(y\) 得反函数
\(y=\frac{2}{1-x}\)，\(x\in\R\)，\(\ x\ne1\).
\end{solution}
\end{problem}

\begin{problem}
设 \(M\)，\(N\) 是直角梯形 \(ABCD\) 两腰的中点，\(DE \perp AB\) 于 \(E\)（如图）. 现将 \(\triangle ADE\) 沿 \(DE\) 折起，使二面角 \(A-DE-B\) 为 \(45^{\circ}\)，此时点 \(A\) 在平面 \(BCDE\) 内的射影恰为点 \(B\)，则 \(M\)，\(N\) 的连线与 \(AE\) 所成角的大小等于\fillinblank{}.

\FigureLayoutDeclare{img/2005/zhejiang_science/q12_fig1.png}{11.0cm}{31bp 24bp 70bp 0bp}
\bitmapfigure[max width=0.74\linewidth]{img/2005/zhejiang_science/q12_fig1.png}

\begin{answer}
\(90^{\circ}\)
\end{answer}

\begin{solution}
建立空间直角坐标系，使平面 \(BCDE\) 为 \(z=0\) 平面，\(EB\) 为 \(x\) 轴正向，\(ED\) 为 \(y\) 轴正向. 设 \(EB=u>0\)，\(ED=v>0\)，则
\(E=(0,0,0)\)，\(B=(u,0,0)\)，\(D=(0,v,0)\)，\(C=(u,v,0)\).
点 \(A\) 在平面 \(BCDE\) 内的射影为 \(B\)，故可设 \(A=(u,0,h)\). 在垂直于 \(DE\) 的截面内，二面角 \(A-DE-B\) 是射线 \(EA\) 与 \(EB\) 的夹角. 该角为 \(45^{\circ}\)，所以 \(h=u\)，即 \(A=(u,0,u)\).

由中点坐标公式，
\(M=\left(\frac{u}{2},\frac{v}{2},\frac{u}{2}\right)\)，\(N=\left(u,\frac{v}{2},0\right)\).
于是
\(\overrightarrow{MN}=\left(\frac{u}{2},0,-\frac{u}{2}\right)\)，\(\overrightarrow{AE}=(-u,0,-u)\)
且
\[
\overrightarrow{MN}\cdot\overrightarrow{AE}
=-\frac{u^2}{2}+\frac{u^2}{2}=0
\]
所以 \(MN\perp AE\)，所成角为 \(90^{\circ}\).
\end{solution}
\end{problem}

\begin{problem}
过双曲线 \( \frac{x^{2}}{a^{2}} - \frac{y^{2}}{b^{2}} = 1 \)（\(a > 0\)，\(b > 0\)）的左焦点且垂直于 \(x\) 轴的直线与双曲线相交于 \(M\)，\(N\) 两点，以 \(MN\) 为直径的圆恰好过双曲线的右顶点，则双曲线的离心率等于\fillinblank{}.

\begin{answer}
\(2\)
\end{answer}

\begin{solution}
设双曲线的半焦距为 \(c\)，则 \(c^2=a^2+b^2\)，左焦点为 \((-c,0)\). 直线 \(x=-c\) 与双曲线相交时，代入方程得
\(\frac{c^2}{a^2}-\frac{y^2}{b^2}=1\)，\(\Longrightarrow\)，\(y^2=\frac{b^4}{a^2}\).
所以 \(M\)，\(N\) 的纵坐标分别为 \(\frac{b^2}{a}\) 和 \(-\frac{b^2}{a}\)，以 \(MN\) 为直径的圆的圆心为 \((-c,0)\)，半径为 \(\frac{b^2}{a}\).

双曲线右顶点为 \((a,0)\)，由题意
\[
a+c=\frac{b^2}{a}=\frac{c^2-a^2}{a}=\frac{(c-a)(c+a)}{a}
\]
因 \(c+a>0\)，约去 \(c+a\) 得 \(c-a=a\)，即 \(c=2a\). 因而离心率
\[
e=\frac{c}{a}=2
\]
\end{solution}
\end{problem}

\begin{problem}
从集合 \(\{O, P, Q, R, S\}\) 与 \(\{0, 1, 2, 3, 4, 5, 6, 7, 8, 9\}\) 中各任取 \(2\) 个元素排成一排（字母和数字均不能重复）. 每排中字母 \(O\)，\(Q\) 和数字 \(0\) 至多只能出现一个的不同排法种数是\fillinblank{}. （用数字作答）

\begin{answer}
\(8424\)
\end{answer}

\begin{solution}
先任取两个字母、两个数字，再将四个元素排列，总排法数为
\[
\mathrm C_5^2\mathrm C_{10}^2 4!=10\times45\times24=10800
\]
不合条件的排法分三类且互不相交：
\begin{enumerate}
\item 同时选中字母 \(O\)，\(Q\). 两个数字任取，有 \(\mathrm C_{10}^2\) 种选法，排法数为 \(\mathrm C_{10}^2 4!=1080\)；
\item 选中 \(O\) 和数字 \(0\)，但未选中 \(Q\). 另一个字母有 \(3\) 种，另一个数字有 \(9\) 种，排法数为 \(3\times9\times4!=648\)；
\item 选中 \(Q\) 和数字 \(0\)，但未选中 \(O\). 同理，排法数为 \(3\times9\times4!=648\).
\end{enumerate}
因此符合条件的排法数为
\[
10800-(1080+648+648)=8424
\]
\end{solution}
\end{problem}

\section{解答题}
\begin{problem}
已知函数 \( f(x) = -\sqrt{3}\sin^{2}x + \sin x\cos x \).
\begin{enumerate}
\item 求 \( f\left(\frac{25\pi}{6}\right) \) 的值；
\item 设 \(\alpha \in (0, \pi)\)，\(f\left(\frac{\alpha}{2}\right) = \frac{1}{4} - \frac{\sqrt{3}}{2}\)，求 \(\sin \alpha\) 的值.
\end{enumerate}

\begin{solution}
\begin{enumerate}
\item 因为 \(\frac{25\pi}{6}=4\pi+\frac{\pi}{6}\)，所以
\(\sin\frac{25\pi}{6}=\frac{1}{2}\)，\(\cos\frac{25\pi}{6}=\frac{\sqrt{3}}{2}\).
因此
\[
f\left(\frac{25\pi}{6}\right)
=-\sqrt{3}\left(\frac{1}{2}\right)^2+\frac{1}{2}\cdot\frac{\sqrt{3}}{2}=0
\]

\item 令 \(\theta=\frac{\alpha}{2}\)，则 \(\theta\in\left(0,\frac{\pi}{2}\right)\). 利用半角公式，原条件化为
\[
-\sqrt{3}\frac{1-\cos\alpha}{2}+\frac{\sin\alpha}{2}
=\frac{1}{4}-\frac{\sqrt{3}}{2}
\]
整理得
\[
\sin\alpha+\sqrt{3}\cos\alpha=\frac{1}{2}
\]
即
\[
2\sin\left(\alpha+\frac{\pi}{3}\right)=\frac{1}{2}
\]
令 \(\delta=\arcsin\frac{1}{4}\)，则 \(0<\delta<\frac{\pi}{6}\). 由于
\(\alpha+\frac{\pi}{3}\in\left(\frac{\pi}{3},\frac{4\pi}{3}\right)\)，只有
\[
\alpha+\frac{\pi}{3}=\pi-\delta
\]
满足条件，所以 \(\alpha=\frac{2\pi}{3}-\delta\). 又
\(\sin\delta=\frac14\)，\(\cos\delta=\frac{\sqrt{15}}4\)，从而
\[
\sin\alpha
=\sin\left(\frac{2\pi}{3}-\delta\right)
=\frac{\sqrt{3}}{2}\cdot\frac{\sqrt{15}}{4}
+\frac{1}{2}\cdot\frac{1}{4}
=\frac{1+3\sqrt{5}}{8}
\]
\end{enumerate}
\end{solution}
\end{problem}

\begin{problem}
已知函数 \( f(x) \) 和 \( g(x) \) 的图象关于原点对称，且 \( f(x) = x^{2} + 2x \).
\begin{enumerate}
\item 求函数 \( g(x) \) 的解析式；
\item 解不等式 \( g(x) \geqslant f(x) - |x - 1| \).
\end{enumerate}

\begin{solution}
\begin{enumerate}
\item 图象关于原点对称意味着点 \((x,f(x))\) 关于原点的对称点 \((-x,-f(x))\) 在函数 \(g\) 的图象上. 因而
\[
g(x)=-f(-x)=-\left[(-x)^2+2(-x)\right]=-x^2+2x
\]

\item 将 \(f\)，\(g\) 的解析式代入不等式，得
\[
-x^2+2x\geqslant x^2+2x-|x-1|
\quad\Longleftrightarrow\quad
|x-1|\geqslant2x^2
\]
当 \(x\geqslant1\) 时，\(|x-1|=x-1\)，不等式化为
\[
2x^2-x+1\leqslant0
\]
其判别式为 \(1-8<0\)，故此时无解.

当 \(x<1\) 时，\(|x-1|=1-x\)，不等式化为
\[
2x^2+x-1\leqslant0
\quad\Longleftrightarrow\quad
(2x-1)(x+1)\leqslant0
\]
所以 \(-1\leqslant x\leqslant\frac12\). 综合两种情况，解集为
\(\left[-1,\frac12\right]\).
\end{enumerate}
\end{solution}
\end{problem}

\begin{problem}
如图，已知椭圆的中心在坐标原点，焦点 \(F_{1}\)，\(F_{2}\) 在 \(x\) 轴上，长轴 \(A_{1}A_{2}\) 的长为 \(4\)，左准线 \(l\) 与 \(x\) 轴的交点为 \(M\)，\(|MA_{1}|: |A_{1}F_{1}| = 2:1\).
\begin{enumerate}
\item 求椭圆的方程；
\item 若直线 \(l_{1}: x = m\)（\(|m| > 1\)），\(P\) 为 \(l_{1}\) 上的动点，使 \(\angle F_{1}PF_{2}\) 最大的点 \(P\) 记为 \(Q\)，求点 \(Q\) 的坐标.（用 \(m\) 表示）
\end{enumerate}

\bitmapfigure[max width=0.60\linewidth]{img/2005/zhejiang_science/q17_fig1.png}

\begin{solution}
\begin{enumerate}
\item 设椭圆的长半轴为 \(a\)，半焦距为 \(c\). 由长轴长为 \(4\)，得 \(a=2\). 左顶点为 \(A_1=(-2,0)\)，左焦点为 \(F_1=(-c,0)\)，左准线为 \(x=-\frac{a^2}{c}=-\frac4c\). 因为 \(M\) 在左准线上，且 \(M\) 位于 \(A_1\) 左侧，所以
\(|MA_1|=\frac4c-2\)，\(|A_1F_1|=2-c\).
由比例条件
\[
\frac4c-2=2(2-c)
\]
得
\(2c^2-6c+4=0\)，\(\Longrightarrow\)，\((c-1)(c-2)=0\).
因椭圆要求 \(0<c<a=2\)，故 \(c=1\). 于是
\(b^2=a^2-c^2=3\)，椭圆方程为
\[
\frac{x^2}{4}+\frac{y^2}{3}=1
\]

\item 由上问，\(F_1=(-1,0)\)，\(F_2=(1,0)\). 设 \(P=(m,y)\)，令 \(u=|y|\). 从 \(P\) 指向两焦点的两个向量的内积为
\[
(m^2-1)+y^2=m^2+u^2-1
\]
其行列式的绝对值为 \(2u\). 因为 \(|m|>1\)，所成角 \(\theta=\angle F_1PF_2\) 为锐角，故
\[
\tan\theta=\frac{2u}{m^2+u^2-1}
\]
对 \(u\geqslant0\) 求导，得
\[
\left(\frac{2u}{m^2+u^2-1}\right)'
=\frac{2(m^2-1-u^2)}{(m^2+u^2-1)^2}
\]
所以当 \(u=\sqrt{m^2-1}\) 时该比值取得最大值，角 \(\theta\) 也取得最大值. 因此
\[
Q\left(m,\pm\sqrt{m^2-1}\right)
\]
\end{enumerate}
\end{solution}
\end{problem}

\begin{problem}
如图，在三棱锥 \(P - ABC\) 中，\(AB \perp BC\)，\(AB = BC = kPA\)，点 \(O\)，\(D\) 分别是 \(AC\)，\(PC\) 的中点，\(OP \perp\) 底面 \(ABC\).
\begin{enumerate}
\item 求证： \(OD \parallel\) 平面 \(PAB\)；
\item 当 \(k = \frac{1}{2}\) 时，求直线 \(PA\) 与平面 \(PBC\) 所成角的大小；
\item 当 \(k\) 取何值时，\(O\) 在平面 \(PBC\) 内的射影恰好为 \(\triangle PBC\) 的重心？
\end{enumerate}

\bitmapfigure[max width=0.74\linewidth]{img/2005/zhejiang_science/q18_fig1.png}

\begin{solution}
\begin{enumerate}
\item 在三角形 \(APC\) 中，\(O\)，\(D\) 分别是 \(AC\)，\(PC\) 的中点. 由三角形中位线定理，
\[
OD\parallel AP
\]
又 \(AP\subset\) 平面 \(PAB\)，所以 \(OD\parallel\) 平面 \(PAB\).

\item 设 \(PA=p>0\)，令 \(s=AB=BC=kp\). 建立空间直角坐标系，取
\(B=(0,0,0)\)，\(A=(s,0,0)\)，\(C=(0,s,0)\).
由于 \(OP\perp\) 底面且 \(O\) 为 \(AC\) 的中点，设
\[
P=\left(\frac{s}{2},\frac{s}{2},h\right)
\]
由 \(|PA|=p\) 得
\[
h^2+\frac{s^2}{2}=p^2
\]
平面 \(PBC\) 的一个法向量为
\[
\bs n=\overrightarrow{BC}\times\overrightarrow{BP}
=\left(sh,0,-\frac{s^2}{2}\right)
\]
取直线 \(PA\) 的方向向量 \(\bs v=\overrightarrow{AP}=\left(-\frac{s}{2},\frac{s}{2},h\right)\)，则
\(|\bs v|=p\)，\(|\bs n|=s\sqrt{h^2+\frac{s^2}{4}}\)，\(|\bs v\cdot\bs n|=s^2h\).
若直线 \(PA\) 与平面 \(PBC\) 所成的角为 \(\theta\)，则
\[
\sin\theta=\frac{|\bs v\cdot\bs n|}{|\bs v|\,|\bs n|}
=\frac{sh}{p\sqrt{h^2+s^2/4}}
\]
当 \(k=\frac12\) 时，\(s=\frac p2\)，且
\(h^2=p^2-\frac{s^2}{2}=\frac{7p^2}{8}\)，\(h=\frac{\sqrt{14}}{4}p\).
代入得
\[
\sin\theta=\frac{\sqrt{14}}{2\sqrt{15}}=\frac{\sqrt{210}}{30}
\]
所以所成角为
\(\displaystyle\theta=\arcsin\frac{\sqrt{210}}{30}\).

\item 由上面的坐标设定，三角形 \(PBC\) 的重心为
\[
G=\frac{P+B+C}{3}=\left(\frac{s}{6},\frac{s}{2},\frac{h}{3}\right)
\]
而
\(O=\left(\frac{s}{2},\frac{s}{2},0\right)\)，\(\overrightarrow{GO}=\left(\frac{s}{3},0,-\frac{h}{3}\right)\).
仍取平面 \(PBC\) 的法向量 \(\bs n=\left(sh,0,-\frac{s^2}{2}\right)\). 点 \(O\) 在平面 \(PBC\) 内的射影为 \(G\) 的充要条件是 \(GO\perp\) 平面 \(PBC\)，即 \(\overrightarrow{GO}\parallel\bs n\). 因此
\(\left(s,0,-h\right)\parallel\left(h,0,-\frac{s}{2}\right)\)，\(\Longrightarrow\)，\(h^2=\frac{s^2}{2}\).
结合 \(h^2+s^2/2=p^2\)，得 \(s^2=p^2\). 因 \(s=kp>0\)，所以 \(k=1\). 反之，\(k=1\) 时上述平行关系成立，故此时射影确为重心.
\end{enumerate}
\end{solution}
\end{problem}

\begin{problem}
袋子 \(A\) 和 \(B\) 中装有若干个均匀的红球和白球，从 \(A\) 中摸出一个红球的概率是 \(\frac{1}{3}\)，从 \(B\) 中摸出一个红球的概率为 \(p\).
\begin{enumerate}
\item 从 \(A\) 中有放回地摸球，每次摸出一个，有 \(3\) 次摸到红球即停止.
\begin{enumerate}
\item 求恰好摸 \(5\) 次停止的概率；
\item 记 \(5\) 次之内（含 \(5\) 次）摸到红球的次数为 \(\xi\)，求随机变量 \(\xi\) 的分布列及数学期望 \(E(\xi)\)；
\end{enumerate}
\item 若 \(A\)，\(B\) 两个袋子中的球数之比为 \(1:2\)，将 \(A\)，\(B\) 中的球装在一起后，从中摸出一个红球的概率是 \(\frac{2}{5}\)，求 \(p\) 的值.
\end{enumerate}

\begin{solution}
\begin{enumerate}
\item 设一次摸到红球和白球的概率分别为 \(r=\frac13\)，\(w=\frac23\).
\begin{enumerate}
\item 恰好第 \(5\) 次停止，等价于前 \(4\) 次恰有 \(2\) 次摸到红球且第 \(5\) 次摸到红球. 因此
\[
P=\mathrm C_4^2r^2w^2r
=6\left(\frac13\right)^3\left(\frac23\right)^2
=\frac{8}{81}
\]

\item 当 \(\xi=0\)，\(1\)，\(2\) 时，在前 \(5\) 次中分别恰有 \(0\)，\(1\)，\(2\) 次红球，故
\[
P(\xi=0)=w^5=\frac{32}{243}
\]
\[
P(\xi=1)=\mathrm C_5^1rw^4=\frac{80}{243}
\]
\[
P(\xi=2)=\mathrm C_5^2r^2w^3=\frac{80}{243}
\]
当 \(\xi=3\) 时，第三次红球在第 \(3\)，\(4\) 或 \(5\) 次出现，所以
\[
P(\xi=3)=r^3+\mathrm C_3^2r^3w+\mathrm C_4^2r^3w^2
=\frac{1}{27}+\frac{2}{27}+\frac{8}{81}
=\frac{17}{81}=\frac{51}{243}
\]
所以分布列为
\[
\begin{tabular}{c|cccc}
\(\xi\) & \(0\) & \(1\) & \(2\) & \(3\) \\
\hline
P & \(\dfrac{32}{243}\) & \(\dfrac{80}{243}\) & \(\dfrac{80}{243}\) & \(\dfrac{51}{243}\)
\end{tabular}
\]
且
\[
E(\xi)=0\cdot\frac{32}{243}+1\cdot\frac{80}{243}+2\cdot\frac{80}{243}+3\cdot\frac{51}{243}=\frac{131}{81}
\]
\end{enumerate}

\item 设袋子 \(A\) 中有 \(N\) 个球，则袋子 \(B\) 中有 \(2N\) 个球. 两袋合并后红球总数为
\(\frac N3+2Np\)，球总数为 \(3N\). 由题意
\(\frac{\frac N3+2Np}{3N}=\frac25\)，\(\Longrightarrow\)，\(\frac19+\frac{2p}{3}=\frac25\).
解得 \(p=\frac{13}{30}\).
\end{enumerate}
\end{solution}
\end{problem}

\begin{problem}
设点 \(A_{n}(x_{n},0)\)，\(P_{n}(x_{n},2^{n - 1})\) 和抛物线 \(C_n: y = x^{2} + a_n x + b_n\)（\(n \in \N^{*}\)），其中 \(a_{n} = -2 - 4n - \frac{1}{2^{n-1}}\)，\(x_{n}\) 由以下方法得到：\(x_{1} = 1\)，点 \(P_{2}(x_{2},2)\) 在抛物线 \(C_{1}: y = x^{2} + a_{1}x + b_{1}\) 上，点 \(A_{1}(x_{1},0)\) 到 \(P_{2}\) 的距离是 \(A_{1}\) 到 \(C_{1}\) 上点的最短距离. \(\cdots\)，点 \(P_{n + 1}(x_{n + 1},2^{n})\) 在抛物线 \(C_{n}: y = x^{2} + a_{n}x + b_{n}\) 上，点 \(A_{n}(x_{n},0)\) 到 \(P_{n + 1}\) 的距离是 \(A_{n}\) 到 \(C_{n}\) 上点的最短距离.
\begin{enumerate}
\item 求 \(x_{2}\) 及 \(C_{1}\) 的方程.
\item 证明 \(\{x_{n}\}\) 是等差数列.
\end{enumerate}

\begin{solution}
\begin{enumerate}
\item 由 \(a_1=-2-4-1=-7\)，知 \(C_1:y=x^2-7x+b_1\). 抛物线在点 \(P_2=(x_2,2)\) 处的切线方向向量为 \((1,2x_2-7)\)，而
\(\overrightarrow{A_1P_2}=(x_2-1,2)\). 最短距离条件说明这两条方向相互垂直，因此
\((x_2-1,2)\cdot(1,2x_2-7)=0\)，\(\Longrightarrow\)，\(x_2-1+2(2x_2-7)=0\).
解得 \(x_2=3\). 再由 \(P_2\in C_1\)，得
\[
2=3^2-7\cdot3+b_1
\]
所以 \(b_1=14\)，从而
\[
C_1:y=x^2-7x+14
\]

\item 对任意 \(n\in\N^*\)，抛物线 \(C_n\) 在点 \(P_{n+1}\) 处的切线方向向量为
\(\left(1,2x_{n+1}+a_n\right)\)，而
\[
\overrightarrow{A_nP_{n+1}}=\left(x_{n+1}-x_n,2^n\right)
\]
由最短距离条件，二者垂直，故
\[
x_{n+1}-x_n+2^n(2x_{n+1}+a_n)=0
\]
假设 \(x_n=2n-1\)，代入 \(a_n=-2-4n-2^{1-n}\)，得
\[
\left(1+2^{n+1}\right)x_{n+1}
=2n-1+2^n\left(2+4n+2^{1-n}\right)
=(2n+1)\left(1+2^{n+1}\right)
\]
因此 \(x_{n+1}=2n+1\). 又 \(x_1=1=2\times1-1\)，由数学归纳法
\(x_n=2n-1\)，\((n\in\N^*)\).
所以 \(x_{n+1}-x_n=2\)，即 \(\{x_n\}\) 是首项为 \(1\)，公差为 \(2\) 的等差数列.
\end{enumerate}
\end{solution}
\end{problem}
